\section{Introduction}\label{sec:introduction}

Can international compatibility fail to emerge even when
each country would gain from a common technical standard?
In industries where interoperability links consumers or
networks across products, a government may expect
compatibility to raise consumer welfare. The same
commitment can change firms' incentives to differentiate
the nonstandardized portions of their products, however,
and those decisions are made \emph{after} governments
announce a technical standard. A forecast that treats
other product characteristics as fixed may therefore
give the wrong answer about which international
agreements countries will accept.

We study an intentionally stylized setting in which
two countries first negotiate a technical-compatibility
union and a third may subsequently seek accession.
The existing members must unanimously approve the
third, and the new entrant inherits the original
compatibility depth. Each country values its domestic
consumers together with the global profits of its
home firm. After the government's decision, firms
independently choose costly residual product attributes,
and then compete in quantities across three separate
national markets. In this environment, an interoperable
\emph{interface} need not make final products identical.

The paper constructs a restricted but sharp
equilibrium counterexample. At one exact rational
parameter vector, with the firms freely reoptimizing
their nonstandardized attributes, no bilateral union
at any strictly positive compatibility depth is
acceptable to its prospective members, and the
third country would rather remain outside than join
the corresponding enlarged international union.
Thus the specified bilateral-first procedure implements
no substantive technical agreement. Yet an
\emph{international} compatible allocation at one
positive depth, conditional on firms subsequently
reoptimizing, would improve \emph{every country's}
welfare relative to complete technical fragmentation.
The obstacle is the path of coalition formation and
the outsider's intermediate option, not a negative
aggregate value of technical compatibility.

A controlled benchmark reveals the role of downstream
design. Hold each firm's residual product attributes
fixed at the design it chooses under separate
standards, while leaving the governments' actions,
market competition and once-paid design costs
unchanged. In that \emph{strategy-restricted}
counterfactual game, a positive-depth international
union is now acceptable to the outsider and the
incumbent members and is preferred by the proposer
to any terminal bilateral union. The predicted
institutional regime therefore reverses when
endogenous product attributes are restored.
Our main all-depth inequalities are supported by
analytical continuation arguments and an independently
reconstructed rational-interval certificate; they
do not rest merely on a numerical policy grid.

The mechanism is not two independently novel
sign-flip results for the incumbent and the outsider.
At the symmetric incompatibility point their marginal
changes are linked by a simple graph-symmetry identity.
Rather, policy-dependent private residual design
changes product-market competition and the national
valuation of exported profits, while the
bilateral-first participation institution maps
the resulting country payoffs into a blocked
multilateral agreement. The government can experience
a positive direct benefit from compatibility yet a
negative \emph{total} derivative once firms adjust
their residual characteristics.

The distinction is important relative to several
prior theories. Network benefits versus heightened
competition and standards-coalition formation are
already familiar \citep{EconomidesFlyer1998,
EconomidesSkrzypacz2003}. Three-country standard
recognition and regional unions were studied by
\citet{GandalShy2001}, with exogenous Salop variety
and subsequent Bertrand pricing.
\citet{CabralSalant2014} show that standardization
can weaken firms' innovation incentives,
and \citet{BuchholzHattori2024} identify a
technology--coalition paradox in public-good
cooperation. \citet{TakaradaKawabata2025}
examine regional versus multilateral regulatory
mutual recognition and the significance of
exogenous differentiation. We do not claim
priority for these general ingredients.
The narrower contribution is to exhibit and
certify a combination of \emph{governmental
sequencing, post-commitment private design,
and Pareto-beneficial but inaccessible final
compatibility} in an explicit market and
national-welfare model.

Our conclusion is correspondingly limited.
The exact model has symmetric countries,
linear-quadratic conditional demand, and
a specified bilateral-first admission rule.
A separately pre-registered experiment with
a nonlinear mapping from product design to
substitutability reproduces the \emph{existence}
contrast at the same economic parameters,
but not at the same compatibility depth:
its initially specified fixed-depth participation
condition fails. No generic result for all
product-design functions or bargaining
institutions is established. Historical
bilateral technical initiatives such as
GSM development and ChaoJi charging cooperation
motivate the timing of early technical alignment
before wider standardization; they are not
claimed to be empirical realizations of the
model's third-country rejection.

Section~\ref{sec:model} defines the
sequential game and separates standardized
interfaces from residual product attributes.
Section~\ref{sec:equilibrium} derives the
Cournot and design continuations,
Section~\ref{sec:results} establishes
the endogenous and frozen-design
coalition results, and Section~\ref{sec:welfare}
shows the remaining Pareto improvement.
Section~\ref{sec:robustness} gives the
nonlinear diagnostic, followed by the
closest-literature comparison, institutional
limitations and proof appendix.
